\chapter{Conclusiones}
El programa desarrollado demuestra un funcionamiento correcto y eficiente al codificar y decodificar imágenes en escala de grises utilizando un autómata celular basado en la Regla 90 modificada (90R), logrando una reconstrucción fiel de las imágenes originales tras aplicar múltiples iteraciones del algoritmo. La estructura modular del código, el uso adecuado de archivos .npy para almacenar los datos codificados y una interfaz gráfica intuitiva permiten al usuario operar el sistema sin complicaciones técnicas, mientras que la visualización grafica facilita la comparación entre las imágenes originales y procesadas. Además, el manejo de las imagenes codificadas mediantre un archi npy contribuye a una experiencia robusta, aunque se debe destacar que la reversibilidad del proceso depende estrictamente de utilizar el mismo número de iteraciones en ambas etapas, lo cual requiere atención por parte del usuario.
